% Converted from README.md
\paragraph{Primer: Pretraga niski}\label{primer-pretraga-niski}

Program pretražuje nisku \texttt{txt} tražeći podnisku \texttt{pat}. U programu se 10 puta vrši pretraga iste niske i iste podniske koristeći 3 različita algoritma. Pomoću Callgrind-a možemo utvrditi koji od algoritama najviše instrukcija izvršava.

Izvorni kod je u datoteci \texttt{pattern\_search.c}. Primer sadrži naivnu pretragu, Boyer-Moore pretragu sa heuristikom lošeg karaktera i Knuth-Morris-Pratt pretragu. Pošto se isti ulaz obrađuje više puta, Callgrind jasno prikazuje cenu svake implementacije.

Primer se prevodi i profajlira iz direktorijuma primera:

\begin{minted}{bash}
$ make
$ valgrind --tool=callgrind ./pattern_search.out
$ callgrind_annotate callgrind.out.*
\end{minted}

U izveštaju treba porediti broj instrukcija u funkcijama \texttt{search\_naive}, \texttt{search\_boyer} i \texttt{search\_kmp}. Dobar signal za optimizaciju nije samo ukupna cena funkcije, već i broj poziva pomoćnih funkcija, na primer računanje tabele prefiksa kod KMP algoritma ili popunjavanje tabele loših karaktera kod Boyer-Moore algoritma.

Za grafički pregled može se otvoriti isti izveštaj u alatu KCachegrind. Tada je korisno krenuti od \texttt{main} funkcije i pratiti inkluzivnu cenu ka funkcijama za pretragu.
